\documentclass[11pt,a4paper]{article}
\usepackage[T1]{fontenc}
\usepackage[margin=25mm]{geometry}
\usepackage{amsmath,amssymb,amsthm,mathtools,microtype,mathrsfs}
\usepackage[hidelinks]{hyperref}
\allowdisplaybreaks
\setlength{\emergencystretch}{3em}
\DeclareMathOperator{\im}{im}
\DeclareMathOperator{\diag}{diag}
\title{The complete native-to-arithmetic class and its original action response}
\author{Split-Zero mathematical continuation}
\date{19 September 2026}
\begin{document}
\maketitle
\begin{abstract}
The actual DNE numerator section has a residue-dual description that omits
nonzero components of every retained simple eigenclass. We construct its
full low, conductor and native completion through every original minor,
calculate its metric and action with all cross terms, and carry the
complete class through the original arithmetic realization. The complete
centered action has an exact rank-two boundary formula at every original
cutoff. Its two boundary rows give precisely the original EC response
ratios, including their signs and both directed leakages. The parity of
the full physical source evaluates the mixed boundary Gram entry exactly
and identifies the product of its two observation amplifications. The
proof retains the original full root polynomial, measure, minimum,
conductor and projection. Absolute determinant allocations and individual
projected arithmetic signs are not evaluated by these finite identities.
\end{abstract}
The originating programme and its existing public attribution are
retained \cite{human:globalization2026,human:splitzero2026}.
The supplied continuation is \cite{prog:arrival02}; source editions and
complete input hashes are recorded in \texttt{SOURCE\_PINS.json}.
All original provider files remain unchanged. The complete additive
finite calculations are written below; their original theta-source
realization is the pinned FCI/FDC provider, whose original hypotheses
are retained. No new certification of that provider's earlier analytic
estimates is implied by this finite intake.

% Additive complete algebra proof; no preamble or bibliography.
% Parent bibliography keys: prog:DNE, prog:arrival02.
\section{The original residue-dual basis, every conductor chart,
and the exact action-defect rank}

\subsection{Original objects and scope}
This calculation proves equations (3)--(17) of the supplied
native-class continuation \cite{prog:arrival02} in the original DNE
coordinates \cite{prog:DNE}. It also evaluates the action-defect rank
for every selected set and constructs the exact transition from every
original row minor to the consecutive chart. All residue expansions and
matrix eliminations used below are proved here.

Keep the actual simple-quartet data and the actual full conductor:
\begin{equation}\tag{NCB1}\label{ncb:data}
\begin{gathered}
a=k+4,\quad q=(a+1)^2,\quad d=q'=(a-7)^2,\quad
v=\operatorname{ord}_0E_A,\quad h=q-v,\quad m=h-d>0,\\
\beta_{a,u,w}=\frac a2+(2u-a)\delta+i(2w-a)\gamma,
\quad 0\leq u,w\leq a,\quad 0<\delta<\tfrac12,\quad\gamma>2,\\
E_A(z)=\sum_{r,t=0}^{8}a_{rt}e^{\beta_{8,r,t}z},
\qquad \mu_j=E_A^{(j)}(0),\qquad
\mu_0=\cdots=\mu_{v-1}=0,\quad\mu_v\ne0 .
\end{gathered}
\end{equation}
Every actual period, branch, coefficient, and finite guard remains
unchanged. The lower roots are
$\beta'_{\alpha'}=\beta_{a-8,u',w'}$, with their original ordering.
There are $d$ lower roots and $q$ upper roots. All roots within either
grid are distinct: equality of real and imaginary parts forces equality
of both indices. Every upper root is nonzero since its real part is
at least $a(1/2-\delta)>0$. Set
\begin{equation}\tag{NCB2}\label{ncb:chi}
\chi(S)=\prod_{\alpha=1}^{q}(S-\beta_\alpha)
       =\sum_{\ell=0}^{q}c_\ell S^\ell,\quad c_q=1,\quad
E=\mathbb C[S]/(\chi),\quad
\mathbb V_{n,\alpha}=\beta_\alpha^n\quad(0\leq n<q).
\end{equation}
Represent elements of $E$ by their unique remainders of degree below
$q$. For an ordered subset $B$ of coordinate indices, $E_B$ inserts
coordinates in increasing index order; $E_B^{\mathsf T}$ extracts them.
These are ordinary transpose conventions, not Hermitian adjoints.

The literal DNE conductor coefficient map
$C_a^{\mathsf T}:\mathbb C^d\to\mathbb C^q$ is
\begin{equation}\tag{NCB3}\label{ncb:conductor}
\sum_\alpha(C_a^{\mathsf T}\eta)_\alpha e^{\beta_\alpha z}
=E_A(z)\sum_{\alpha'}\eta_{\alpha'}e^{\beta'_{\alpha'}z}.
\end{equation}
Its upper $(u,w)$, lower $(u',w')$ entry is
$a_{u-u',w-w'}$ when both differences lie in $\{0,\ldots,8\}$, and
zero otherwise, because
$\beta_{8,r,t}+\beta_{a-8,u',w'}=\beta_{a,u'+r,w'+t}$.
No divided-native scalar is inserted in this map.

The map $C_a^{\mathsf T}$ is injective. If its value is zero, the
product in \eqref{ncb:conductor} vanishes identically. Since
$\mu_v\ne0$, $E_A$ is nonzero on an open disc; the lower exponential
polynomial vanishes on that disc and hence identically. Its first $d$
derivatives at zero form the lower Vandermonde times $\eta$. The
determinant is
$\prod_{\alpha'<\beta'}(\beta'_{\beta'}-\beta'_{\alpha'})\ne0$,
so $\eta=0$.
To verify this determinant formula, its determinant is an alternating
polynomial in the $d$ root variables, hence divisible by every
$\beta'_{\beta'}-\beta'_{\alpha'}$. Its total degree is
$0+1+\cdots+(d-1)$, the degree of their product. The coefficient of
$(\beta'_1)^0(\beta'_2)^1\cdots(\beta'_d)^{d-1}$ is one in both
polynomials, so their constant quotient is one.

Let $P_v:\mathbb C^h\to\mathbb C^q$ insert $v$ leading zero
coordinates. The full product rule applied to \eqref{ncb:conductor}
gives
\begin{equation}\tag{NCB4}\label{ncb:moment-conductor}
\begin{gathered}
\mathbb VC_a^{\mathsf T}=P_vJ,\qquad
J=P_v^{\mathsf T}\mathbb VC_a^{\mathsf T},\\
J_{i,\alpha'}=\sum_{r=0}^{i}
\binom{v+i}{r}\mu_{v+i-r}(\beta'_{\alpha'})^r
\qquad(0\leq i<h).
\end{gathered}
\end{equation}
All derivatives below order $v$ vanish. At order $v+i$, the product
rule originally ranges over $0\leq r\leq v+i$; terms with $r>i$
vanish because their conductor derivative has order below $v$.
Thus the displayed truncation retains every nonzero conductor
derivative. Since $\mathbb V$ and $P_v$ are injective, $J$ has rank $d$.

\subsection{Complete residue-dual moment basis}
For each upper root let
$p_\alpha(S)=\chi(S)/((S-\beta_\alpha)\chi'(\beta_\alpha))$.
Define the residue isomorphism and its moment-coordinate version:
\begin{equation}\tag{NCB5}\label{ncb:residue}
\mathcal L:\mathbb C^q\longrightarrow E,\qquad
\mathcal L(\nu)=\sum_\alpha\chi'(\beta_\alpha)\nu_\alpha p_\alpha,
\qquad F=\mathcal L\mathbb V^{-1}.
\end{equation}
Evaluation at the roots gives
$(\mathcal L^{-1}f)_\alpha=f(\beta_\alpha)/\chi'(\beta_\alpha)$,
so $\mathcal L$ is invertible. The roots are simple and no derivative
factor vanishes. The formula for a simple-pole residue gives
\[
\sum_\alpha\operatorname{Res}_{S=\beta_\alpha}
\frac{\mathcal L(\nu)(S)P(S)}{\chi(S)}\,dS
=\sum_\alpha\nu_\alpha P(\beta_\alpha).
\]

Define, for $0\leq n<q$,
\begin{equation}\tag{NCB6}\label{ncb:pn}
\mathfrak p_n(S)=\sum_{\ell=n+1}^{q}c_\ell S^{\ell-n-1}.
\end{equation}
Their distinct monic degrees are $q-1-n$, so they form a basis of $E$.
Their complete residue duality is
\begin{equation}\tag{NCB7}\label{ncb:duality}
\sum_\alpha\operatorname{Res}_{S=\beta_\alpha}
\frac{\mathfrak p_n(S)S^j}{\chi(S)}\,dS=\delta_{nj}
\quad(0\leq j,n<q),\qquad Fe_n=\mathfrak p_n .
\end{equation}
To prove the first equality, the sum of finite residues equals the
$S^{-1}$ coefficient at large $S$: partial fractions consist of a
polynomial and terms $b_\alpha/(S-\beta_\alpha)$, and each such term
has both residue and $S^{-1}$ coefficient $b_\alpha$.
For $j<n$, the numerator has degree at most $q-2$, so this coefficient
is zero. For $j=n$, it has degree $q-1$ and leading coefficient one.
For $j>n$, the exact identity
$S^{n+1}\mathfrak p_n=\chi-\sum_{\ell=0}^{n}c_\ell S^\ell$
implies
\[
\frac{\mathfrak p_nS^j}{\chi}
=S^{j-n-1}
-\frac{S^{j-n-1}\sum_{\ell=0}^{n}c_\ell S^\ell}{\chi}.
\]
The first term is a polynomial, and the second numerator has degree
at most $j-1\leq q-2$. Both have zero $S^{-1}$ coefficient.
This proves the first equality. The pairings of $Fe_n$ with $S^j$
are $(\mathbb V\mathbb V^{-1}e_n)_j=\delta_{nj}$, so the proved
basis duality identifies $Fe_n$ with $\mathfrak p_n$.

Fix any actual $d$-row minor
$I\subset\{0,\ldots,h-1\}$ with $\det J_I\ne0$,
where $J_I=E_I^{\mathsf T}J$. Write $J(I)=I^c$ and
\begin{equation}\tag{NCB8}\label{ncb:native}
T_I=\{v+j:j\in J(I)\},\qquad
Z_I=\mathbb V^{-1}P_vE_{J(I)},\qquad
L_I=\mathcal LZ_I=FE_{T_I}.
\end{equation}
Hence $\operatorname{im}L_I$ is exactly the span of
$\{\mathfrak p_n:n\in T_I\}$ in the original selected coordinates.

For every upper root $\lambda$, the finite geometric sums give
\[
\sum_{n=0}^{q-1}\lambda^n\mathfrak p_n(S)
=\sum_{\ell=1}^{q}c_\ell
\frac{S^\ell-\lambda^\ell}{S-\lambda}
=\frac{\chi(S)-\chi(\lambda)}{S-\lambda}.
\]
Consequently the original eigenclass is
\begin{equation}\tag{NCB9}\label{ncb:eigenclass}
v_\lambda=\frac{\chi(S)}{(S-\lambda)\chi'(\lambda)}
=\sum_{n=0}^{q-1}t_n\mathfrak p_n,\qquad
t_n=\frac{\lambda^n}{\chi'(\lambda)},\qquad
v_\lambda\notin\operatorname{im}L_I.
\end{equation}
Every $t_n$ is nonzero because $\lambda\ne0$ and
$\chi'(\lambda)\ne0$. A proper coordinate subset of the displayed
basis cannot contain a vector whose every coordinate is nonzero.
This proves the last assertion for every original admitted minor.

\subsection{Complete low, conductor, and native decomposition}
For an arbitrary $f\in E$, put $t=F^{-1}f$ and
$t_{\geq v}=P_v^{\mathsf T}t$. Define maps with their full types:
\begin{equation}\tag{NCB10}\label{ncb:HQ}
\begin{gathered}
H_I=J_I^{-1}E_I^{\mathsf T}:\mathbb C^h\to\mathbb C^d,\qquad
Q_I=E_{J(I)}^{\mathsf T}-J_{J(I)}H_I:
       \mathbb C^h\to\mathbb C^m,\\
\eta_I=H_It_{\geq v},\qquad z_I=Q_It_{\geq v}.
\end{gathered}
\end{equation}
Checking first the $I$ rows and then their complement proves
\begin{equation}\tag{NCB11}\label{ncb:split}
\begin{gathered}
JH_I+E_{J(I)}Q_I=I_h,\qquad H_IJ=I_d,\qquad Q_IJ=0,\\
H_IE_{J(I)}=0,\qquad Q_IE_{J(I)}=I_m.
\end{gathered}
\end{equation}
Together with $FP_vJ=\mathcal LC_a^{\mathsf T}$ from
\eqref{ncb:moment-conductor}, these give the full identity
\begin{equation}\tag{NCB12}\label{ncb:decomposition}
f=\underbrace{\sum_{n<v}t_n\mathfrak p_n}_{f^{\rm low}}
 +\underbrace{\mathcal L(C_a^{\mathsf T}\eta_I)}_{f_I^{\rm cond}}
 +\underbrace{L_Iz_I}_{f_I^{\rm nat}} .
\end{equation}
These three spaces form a direct sum. The low space has arbitrary
first $v$ moment coordinates and zero remaining coordinates; both
other spaces have their first $v$ coordinates zero. On the remaining
coordinates, \eqref{ncb:split} proves
$\mathbb C^h=\operatorname{im}J\oplus\operatorname{im}E_{J(I)}$.
Thus the full basis map
\begin{equation}\tag{NCB13}\label{ncb:full-chart}
B_I=[\,FE_{\{0,\ldots,v-1\}},\ \mathcal LC_a^{\mathsf T},\ L_I\,]:
\mathbb C^v\oplus\mathbb C^d\oplus\mathbb C^m\longrightarrow E
\end{equation}
is invertible, also when its first block is empty ($v=0$).

For $f=v_\lambda$, all selected coordinates of $t_{\geq v}$ are
nonzero and $d\geq1$, so $\eta_I\ne0$. Injectivity of the full
conductor map and of $\mathcal L$ gives
$v_{\lambda,I}^{\rm cond}\ne0$. If $v>0$ then
$v_\lambda^{\rm low}\ne0$ as well. The direct sum just proved implies
\begin{equation}\tag{NCB14}\label{ncb:nonzero}
v_\lambda^{\rm low}+v_{\lambda,I}^{\rm cond}\ne0 .
\end{equation}
This is an algebraic direct sum, not an orthogonal splitting.

The exact relative quotient represented by the native section is
\begin{equation}\tag{NCB15}\label{ncb:relative}
\frac{E}
 {\operatorname{im}(FE_{\{0,\ldots,v-1\}})
          +\operatorname{im}(\mathcal LC_a^{\mathsf T})}
\ \xrightarrow[\ [f]\mapsto Q_IP_v^{\mathsf T}F^{-1}f\ ]{\ \cong\ }
\mathbb C^m .
\end{equation}
Its inverse is $z\mapsto[L_Iz]$. Its kernel and both inverse identities
follow from \eqref{ncb:split}. This proves the exact map relating the
full arithmetic class and its native representative; it retains their
proved nonzero difference in $E$.

\subsection{Every original minor to every other minor}
For another admitted $d$-row minor $K$, let
\begin{equation}\tag{NCB16}\label{ncb:transition}
D_{K\leftarrow I}=H_KE_{J(I)}:\mathbb C^m\to\mathbb C^d,\qquad
M_{K\leftarrow I}=Q_KE_{J(I)}:\mathbb C^m\to\mathbb C^m .
\end{equation}
Apply \eqref{ncb:split} for $K$ to $E_{J(I)}$ to obtain
\begin{equation}\tag{NCB17}\label{ncb:section-transition}
\begin{gathered}
E_{J(I)}=JD_{K\leftarrow I}
              +E_{J(K)}M_{K\leftarrow I},\\
L_I=\mathcal LC_a^{\mathsf T}D_{K\leftarrow I}
              +L_KM_{K\leftarrow I}.
\end{gathered}
\end{equation}
This is the full changed-representative identity with its conductor
correction. Apply $Q_I$ to its first line to obtain
$M_{I\leftarrow K}M_{K\leftarrow I}=I_m$, and exchange the minors
to obtain the reverse product. Therefore
\begin{equation}\tag{NCB18}\label{ncb:inverse}
M_{K\leftarrow I}^{-1}=M_{I\leftarrow K},\qquad
\eta_K=\eta_I+D_{K\leftarrow I}z_I,\quad
z_K=M_{K\leftarrow I}z_I,\quad
f_K^{\rm low}=f_I^{\rm low}.
\end{equation}
The component identities follow by applying $H_K,Q_K$ to
$t_{\geq v}=J\eta_I+E_{J(I)}z_I$. In particular
\[
f_I^{\rm nat}=f_K^{\rm nat}
       +\mathcal LC_a^{\mathsf T}D_{K\leftarrow I}z_I,\qquad
f_K^{\rm cond}=f_I^{\rm cond}
       +\mathcal LC_a^{\mathsf T}D_{K\leftarrow I}z_I.
\]
Their total is unchanged because the same correction occurs on
opposite sides when the decompositions are compared.

For a third admitted minor $H$, applying $Q_H,H_H$ to
\eqref{ncb:section-transition} proves the oriented composition laws
\begin{equation}\tag{NCB19}\label{ncb:cocycle}
\begin{gathered}
M_{H\leftarrow I}=M_{H\leftarrow K}M_{K\leftarrow I},\\
D_{H\leftarrow I}
 =D_{K\leftarrow I}+D_{H\leftarrow K}M_{K\leftarrow I}.
\end{gathered}
\end{equation}

The exact determinant including ordering signs is also explicit.
For $I=(i_0<\cdots<i_{d-1})$ put
$\epsilon_I=(-1)^{\sum_{r=0}^{d-1}(i_r-r)}$.
The row permutation taking full increasing order to $(I,J(I))$
has this sign, since selected row $i_r$ crosses $i_r-r$ earlier
unselected rows. Under that permutation $W_I=[J,E_{J(I)}]$ is
block lower triangular with diagonal blocks $J_I,I_m$, so
$\det W_I=\epsilon_I\det J_I$. The first line of
\eqref{ncb:section-transition} gives
$W_I=W_K\left[\begin{smallmatrix}I_d&D\\0&M\end{smallmatrix}\right]$.
Consequently
\begin{equation}\tag{NCB20}\label{ncb:dettransition}
\det M_{K\leftarrow I}
 =\frac{\epsilon_I\det J_I}{\epsilon_K\det J_K}.
\end{equation}

The full coordinate transformation, including the unchanged low
coordinates, is
\begin{equation}\tag{NCB21}\label{ncb:full-transition}
T_{K\leftarrow I}=
\begin{pmatrix}I_v&0&0\\0&I_d&D_{K\leftarrow I}\\
0&0&M_{K\leftarrow I}\end{pmatrix},
\qquad B_I=B_KT_{K\leftarrow I}.
\end{equation}
For every original arithmetic Hermitian metric $G_N$ and the full
operator $A$ of multiplication by $S$ modulo $\chi$, set
$G_{I,N}=B_I^*G_NB_I$ and $A_I=B_I^{-1}AB_I$. Substitution yields
\begin{equation}\tag{NCB22}\label{ncb:gram-action}
G_{I,N}=T_{K\leftarrow I}^*G_{K,N}T_{K\leftarrow I},\qquad
A_I=T_{K\leftarrow I}^{-1}A_KT_{K\leftarrow I}.
\end{equation}
All nine Gram blocks, including all six off-diagonal pairings, are
retained. Multiplication of these identities gives
\begin{equation}\tag{NCB23}\label{ncb:centered-action}
\begin{split}
A_I^*G_{I,N}+G_{I,N}A_I-aG_{I,N}
={}&T_{K\leftarrow I}^*
(A_K^*G_{K,N}+G_{K,N}A_K-aG_{K,N})T_{K\leftarrow I}.
\end{split}
\end{equation}
Thus the full arithmetic action is preserved under the original
coordinate transition with its explicit conductor corrections.

\subsection{Consecutive chart with all conductor derivatives}
Let $I_0=\{0,\ldots,d-1\}$. Equation \eqref{ncb:moment-conductor}
gives the exact factorization
\begin{equation}\tag{NCB24}\label{ncb:triangular}
J_{I_0}=\mathsf T\mathbb V',\qquad
\mathsf T_{i,r}=
\begin{cases}\binom{v+i}{r}\mu_{v+i-r},&0\leq r\leq i<d,\\
0,&r>i,\end{cases}
\quad \mathbb V'_{r,\alpha'}=(\beta'_{\alpha'})^r .
\end{equation}
The lower-triangular diagonal entries are
$\binom{v+i}{i}\mu_v$. Taking the determinant of the actual product
therefore proves
\begin{equation}\tag{NCB25}\label{ncb:consecutive-det}
\det J_{I_0}
=\mu_v^d\left(\prod_{i=0}^{d-1}\binom{v+i}{i}\right)
 \prod_{\alpha'<\beta'}(\beta'_{\beta'}-\beta'_{\alpha'})\ne0 .
\end{equation}
All $\mu_j$ with $j>v$ remain in the matrix. Their absence from its
determinant is the proved effect of triangularity.

Its inverse retains every such derivative. For $x\in\mathbb C^d$,
solve $\mathsf T b=x$ by the finite recursion
\begin{equation}\tag{NCB26}\label{ncb:triangular-inverse}
b_0=\frac{x_0}{\mu_v},\qquad
b_i=\frac{x_i-\sum_{r=0}^{i-1}
          \binom{v+i}{r}\mu_{v+i-r}b_r}
          {\binom{v+i}{i}\mu_v}\quad(1\leq i<d).
\end{equation}
Write
$\chi_{\rm low}(S)=\prod_{\alpha'}(S-\beta'_{\alpha'})
=\sum_{\ell=0}^{d}c_\ell^{\rm low}S^\ell$
and
$\mathfrak p_r^{\rm low}(S)
=\sum_{\ell=r+1}^{d}c_\ell^{\rm low}S^{\ell-r-1}$.
The already proved residue-dual basis applied to $\chi_{\rm low}$
evaluates the inverse lower Vandermonde, giving
\begin{equation}\tag{NCB27}\label{ncb:full-inverse}
(J_{I_0}^{-1}x)_{\alpha'}
=\frac{\sum_{r=0}^{d-1}b_r\mathfrak p_r^{\rm low}(\beta'_{\alpha'})}
          {\chi_{\rm low}'(\beta'_{\alpha'})}.
\end{equation}
Indeed the vector on the right has lower moments $b$ by
\eqref{ncb:duality}, and those moments satisfy $\mathsf T b=x$.

For an arbitrary original chart $I$ and its native vector $z$, take
$x=E_{I_0}^{\mathsf T}E_{J(I)}z$ in
\eqref{ncb:triangular-inverse}--\eqref{ncb:full-inverse}. Their result
is exactly $D_{I_0\leftarrow I}z$. The native coordinates in the
consecutive chart are, for $d\leq j<h$,
\begin{equation}\tag{NCB28}\label{ncb:explicit-transition}
\begin{split}
(M_{I_0\leftarrow I}z)_{j-d}
={}&(E_{J(I)}z)_j\\
&-\sum_{\alpha'}\sum_{r=0}^{j}
\binom{v+j}{r}\mu_{v+j-r}(\beta'_{\alpha'})^r
                    (D_{I_0\leftarrow I}z)_{\alpha'} .
\end{split}
\end{equation}
This evaluates the chart transition in original coefficients, roots,
and derivatives. For the full eigenclass one instead sets
$x_i=\lambda^{v+i}/\chi'(\lambda)$ to calculate $\eta_{I_0}$,
and then sets
$z_{I_0,j-d}=t_{v+j}-(J\eta_{I_0})_j$.

The selected consecutive set is
$T_{I_0}=\{v+d,\ldots,q-1\}=\{q-m,\ldots,q-1\}$.
Its residue-dual basis elements have all monic degrees
$m-1,\ldots,0$, so their triangular coefficient matrix proves
\begin{equation}\tag{NCB29}\label{ncb:consecutive-image}
\operatorname{im}L_{I_0}=\mathbb C[S]_{\leq m-1}\subset E .
\end{equation}

\subsection{Action defect and exact rank for every selected set}
Multiplication by $S$ modulo the complete $\chi$ satisfies
\begin{equation}\tag{NCB30}\label{ncb:action}
A\mathfrak p_0=-c_0\mathfrak p_{q-1},\qquad
A\mathfrak p_n=\mathfrak p_{n-1}-c_n\mathfrak p_{q-1}
\quad(1\leq n<q).
\end{equation}
For the first equality use $S\mathfrak p_0=\chi-c_0$;
for the second use the literal polynomial identity
$S\mathfrak p_n=\mathfrak p_{n-1}-c_n$.
Here $\mathfrak p_{q-1}=1$.
For any selected $T\subset\{0,\ldots,q-1\}$, the omitted-coordinate
action is therefore
\begin{equation}\tag{NCB31}\label{ncb:RT}
\begin{gathered}
R_T=E_{T^c}^{\mathsf T}F^{-1}AFE_T,\\
(R_T)_{\ell n}
=\mathbf1_{\{n\geq1,\ \ell=n-1\}}
       -c_n\mathbf1_{\{\ell=q-1\}}
\qquad(\ell\in T^c,\ n\in T).
\end{gathered}
\end{equation}

Define the descending starts and their number by
$B(T)=\{n\in T:n\geq1,\ n-1\notin T\}$ and $b(T)=|B(T)|$.
The exact rank, including nongeneric vanishing coefficients, is
\begin{equation}\tag{NCB32}\label{ncb:rank}
\begin{gathered}
\operatorname{rank}R_T=b(T)+\varepsilon(T),\\
\varepsilon(T)=
\begin{cases}
1,&q-1\notin T\ \text{and some }n\in T\setminus B(T)
                                    \text{ has }c_n\ne0,\\
0,&\text{otherwise}.
\end{cases}
\end{gathered}
\end{equation}
For each $n\in B(T)$, row $n-1$ has exactly a single one in column
$n$. These $b(T)$ rows are independent and none is row $q-1$.
All other rows vanish except possibly row $q-1$, present exactly when
$q-1\notin T$. That row is $(-c_n)_{n\in T}$.
Subtracting its coordinates on columns $B(T)$ using the independent
rows leaves its coordinates on $T\setminus B(T)$; these add one rank
precisely in the case displayed. This proves \eqref{ncb:rank}.

For every nonempty proper $T$, the rank is positive. If $B(T)$ is
nonempty, this follows immediately. Otherwise descending from any
member of $T$ cannot leave $T$ before reaching zero, so $T$ is an
initial interval containing zero. Properness implies $q-1\notin T$,
and $c_0=(-1)^q\prod_\alpha\beta_\alpha\ne0$ gives
$\varepsilon(T)=1$. For the consecutive set $T_{I_0}$ there is just
one descending start $q-m=v+d>0$, and $q-1\in T_{I_0}$.
Thus its rank is exactly one.

\subsection{Original metric residual and its full Gram}
Let $G_N$ be any original positive Hermitian arithmetic quotient metric
on $E$, fix $T=T_I$, and set $L=FE_T$ and $\Gamma=L^*G_NL$.
The metric compression and residual, both with domain $\mathbb C^m$,
are
\begin{equation}\tag{NCB33}\label{ncb:compression}
A_{\rm nat,N}=\Gamma^{-1}L^*G_NAL,\qquad
\mathcal R_N=AL-LA_{\rm nat,N}.
\end{equation}
Write $\widetilde G=F^*G_NF$, and use $C=T^c$ solely as a block
index in the next formulas. Set
\begin{equation}\tag{NCB34}\label{ncb:lift}
\begin{gathered}
W_G=F(E_C-E_T\widetilde G_{TT}^{-1}\widetilde G_{TC}),\\
S_G=\widetilde G_{CC}
 -\widetilde G_{CT}\widetilde G_{TT}^{-1}\widetilde G_{TC}.
\end{gathered}
\end{equation}
This leaves the original $G_N$ unchanged. Direct multiplication gives
$L^*G_NW_G=0$ and $E_C^{\mathsf T}F^{-1}W_G=I_{q-m}$.
Hence $W_G$ is an injective map onto the $G_N$-orthogonal complement
of $\operatorname{im}L$, whose dimension is $q-m$; its inverse there
is literal omitted-coordinate extraction.
By \eqref{ncb:compression}, $\mathcal R_N$ has image in that
complement and omitted coordinates $R_T$. Uniqueness of this exact
lift proves
\begin{equation}\tag{NCB35}\label{ncb:residual}
\mathcal R_N=W_GR_T,\qquad
\mathcal R_N^*G_N\mathcal R_N=R_T^*S_GR_T,\qquad
\operatorname{rank}\mathcal R_N=\operatorname{rank}R_T .
\end{equation}
For the Gram identity, multiply \eqref{ncb:lift} to obtain
$W_G^*G_NW_G=S_G$, retaining its cross terms.
It is positive definite since $W_G$ is injective and $G_N$ is
positive definite. This proves the residual rank at every original
cutoff, with its exact value \eqref{ncb:rank}, and proves nonvanishing
for every original nonempty proper chart. No size estimate follows
from rank alone.

\subsection{What this finite calculation establishes}
The residue-dual basis, numerator image, every eigenclass
decomposition, nonzero omitted component, original companion action,
and complete consecutive determinant prove equations (3)--(17) of
\cite{prog:arrival02}. Equations
\eqref{ncb:transition}--\eqref{ncb:centered-action} provide the exact
comparison between every original minor and the consecutive chart.
The derivative recursion \eqref{ncb:triangular-inverse} and
\eqref{ncb:full-inverse} evaluate this comparison without deleting
conductor derivatives. The rank formula \eqref{ncb:rank} and residual
Gram \eqref{ncb:residual} strengthen the finite action calculation.

Under the programme's injective arithmetic realization, the nonzero
omitted vector \eqref{ncb:nonzero} remains nonzero. Its map from the
full remainder space to the native relative quotient is precisely
\eqref{ncb:relative}; changing the section is precisely
\eqref{ncb:section-transition}. The arithmetic pairing of the full
class consequently uses the complete block congruence
\eqref{ncb:gram-action} and \eqref{ncb:centered-action}.
This calculation does not identify the native section with the
original observation kernel or evaluate an absolute allocation or
individual projected sign.


% Complete actual-class completion, 19 September 2026.
% Original DNE and FCI sources are unchanged.
\section{The full arithmetic class, its joint Gram, and its action}

Retain the simple quartet, the original period and branch, the full
nonzero conductor $E_A$, and the original full arithmetic unit.
Write $a=k+4$, $k=4\ell+1\ge9$,
$q=(a+1)^2$, $q'=(a-7)^2$, $v=\operatorname{ord}_0E_A\le80$,
$m=q-q'-v$, $c=a/2$, and
\[
 \beta_{a,r,t}=c+(2r-a)\delta+i(2t-a)\gamma,\qquad
 0<\delta<\tfrac12,\quad \gamma>2,\qquad
 \chi(S)=\prod_{r,t=0}^{a}(S-\beta_{a,r,t}).
 \tag{NCRP1}
\]
In particular $q'>0$, $m>0$, every root is simple, and
$\Re\beta\ge a(1/2-\delta)>0$. No root or conductor factor is
removed. The residue-dual polynomial frame and original DNE minor
are those proved in the preceding section. Let
$\mathscr P:\mathbb C^q\to E=\mathbb C[S]/(\chi)$ send the standard
column $e_n$ to $\mathfrak p_n$.
Thus $\mathscr P=\mathcal L\mathbb V^{-1}$, exactly in FCI33--34,
and multiplication by $S$ has matrix $C$ in this frame:
\[
 C_{\ell n}=\mathbf1_{\{n\ge1,\ell=n-1\}}
              -c_n\mathbf1_{\{\ell=q-1\}},\qquad
 \chi(S)=\sum_{n=0}^{q}c_nS^n,\quad c_q=1.
 \tag{NCRP2}
\]
This use of $C$ denotes the full packet action matrix; the original
conductor coefficient matrix continues to be written $C_a$.

\subsection{A square completion using the actual minor}

Let $E_{<v}:\mathbb C^v\to\mathbb C^q$ insert the first $v$
coordinates and retain the original $P_v,J_v,\mathcal I,\mathcal J$.
Define the square coefficient map
\[
 K_{\mathcal I}=
 [E_{<v},\,P_vJ_v,\,P_vE_{\mathcal J}],\qquad
 F_{\mathcal I}=\mathscr P K_{\mathcal I}=[F_0,F_C,L_a].
 \tag{NCRP3}
\]
Here $F_0$ has $v$ columns, $F_C=\mathcal L C_a^{\mathsf T}$
has $q'$ columns, and $L_a=\mathcal L Z_{\exp}$ has $m$ columns.
All three maps have codomain the original polynomial packet $E$.
For a moment column $t=(t_{<v},t_{\ge v})$ the inverse is explicitly
\[
 K_{\mathcal I}^{-1}t=
 \begin{pmatrix}
 t_{<v}\\
 (J_v)_{\mathcal I}^{-1}(t_{\ge v})_{\mathcal I}\\
 (t_{\ge v})_{\mathcal J}
       -(J_v)_{\mathcal J}(J_v)_{\mathcal I}^{-1}
                                (t_{\ge v})_{\mathcal I}
 \end{pmatrix}.
 \tag{NCRP4}
\]
Indeed substitution in the $\mathcal I$ and $\mathcal J$ rows,
and then in the first $v$ rows, gives exactly $t$.
Thus $K_{\mathcal I}$ and $F_{\mathcal I}$ are isomorphisms.
This proves a full arithmetic representative map, not only a
relative quotient map of dimension $m$.

For a retained simple eigenclass $v_\lambda$ let
$t_{\lambda,n}=\lambda^n/\chi'(\lambda)$ and
\[
 w_\lambda=(u_0,\eta,z)^{\mathsf T}
                  =K_{\mathcal I}^{-1}t_\lambda .
 \tag{NCRP5}
\]
The preceding residue calculation gives the exact equality
$v_\lambda=F_0u_0+F_C\eta+L_az$.
The vector $\eta$ is nonzero, because $q'>0$, all entries of
$(t_{\lambda,\ge v})_{\mathcal I}$ are nonzero, and the retained
minor is invertible. If $v>0$, $u_0$ is also nonzero.
Since $F_{\mathcal I}$ is injective, the omitted vector
$F_0u_0+F_C\eta$ is nonzero. These assertions hold for the actual
minor selected in DNE; they do not replace it by the consecutive chart.

\subsection{The original full-unit arithmetic injection}

The map into theta cohomology retains the source $F_h$ of FCI35:
$\mathcal M F_h=v_h=2\xi/h$ and $h(D)F_h=\Theta\phi_*$.
On the simple-quartet branch,
$h(s)=\prod_{\alpha\in\{\rho,\bar\rho,1-\rho,1-\bar\rho\}}(s-\alpha)$.
For $P\in E$ set, with the original Euler variables,
\[
 \iota_a(P)=
 [r_\chi(P)(D_1+\cdots+D_a)F_h^{\otimes a}]
            \in Q_\theta^{\otimes a}.
 \tag{NCRP6}
\]
Here $r_\chi$ is the increasing-degree representative below $q$.
We recall the complete finite reason that this packet map is injective.
In the tensor algebra
$\mathbb C[s_1,\ldots,s_a]/(h(s_1),\ldots,h(s_a))$,
evaluation at all ordered root tuples is an isomorphism: apply the
four-root Vandermonde separately in every coordinate.
Every pair of counts $(r,t)$, $0\le r,t\le a$, occurs among such
tuples, since the real-sign and imaginary-sign choices can be made
independently in each tensor leg. The resulting distinct sums are
exactly the roots in NCRP1. Hence $P(s_1+\cdots+s_a)$ vanishes
in that tensor algebra precisely when $\chi$ divides $P$.
The cyclic annihilator is therefore the full $\chi$.

The original Taylor map on the theta quotient sends NCRP6 to
$\upsilon_h^{\otimes a}P(A_h^{(a)})1$.
Every component of the original primary unit $\upsilon_h=j_h(v_h)$
is nonzero, so multiplication by it is invertible. It cannot turn
a nonzero packet class into zero.
For completeness, multiples of $\chi(s_1+\cdots+s_a)$ lie in the
ideal $(h(s_1),\ldots,h(s_a))$ by the preceding tensor evaluation.
Successive monic division in $s_1,\ldots,s_a$ writes such a polynomial
as $\sum_i h(s_i)R_i$. Substitution of the original differential
operators supplies its theta primitive, leg by leg, with the original
$(-1)^{i-1}$ tensor-differential signs. Thus the representative
independence and the Taylor injectivity apply to the same map.
This is the FCI35/FDC construction \cite{prog:FCI}, with its full
unit and all ordered tuples retained.

Consequently
\[
 \iota_a(v_\lambda)
 =\iota_a(F_0u_0)+\iota_a(F_C\eta)+\iota_a(L_az),
 \qquad
 \iota_a(F_0u_0+F_C\eta)\ne0.
 \tag{NCRP7}
\]
Taking four independent copies of these maps preserves every
equation. The omitted summands cannot be removed as theta boundaries.

\subsection{The full original arithmetic metric}

Let $dm_a^{\rm ar}(y)=m_a^{\rm ar}(y)\,dy$,
$m_a^{\rm ar}=w_h^{*a}$, and
$w_h(y)=|v_h(1/2+iy)|^2/(2\pi)$.
The source has its complete mass $(\int_{\mathbb R}w_h)^a$.
Let $H_N^{\rm ar}$ be its moment Gram in the increasing powers of
$S=c+iy$. For the literal remainder map
$J_N^{\chi}:\mathbb C[S]_{\le N}\to E$ and $N\ge q-1$, set
\[
 G_N=(J_N^\chi(H_N^{\rm ar})^{-1}(J_N^\chi)^*)^{-1},
 \qquad
 T_N=(H_N^{\rm ar})^{-1}(J_N^\chi)^*G_N.
 \tag{NCRP8}
\]
The norm of a class is the minimum over its full relation fibre.
To verify the formula, $J_N^\chi T_N=I$ and
$(T_Nx)^*H_N^{\rm ar}r=0$ whenever $J_N^\chi r=0$.
Every other lift is $T_Nx+r$, so its squared norm is
$x^*G_Nx+r^*H_N^{\rm ar}r$. This proves both the minimum and
$T_N^*H_N^{\rm ar}T_N=G_N$ with the whole relation space
$\chi\mathbb C[S]_{\le N-q}$ retained. It is the classical positive
Schur calculation \cite{human:boyd-vandenberghe}, proved here in
its complex Hermitian coordinates.

All polynomial moments are finite for the admitted arithmetic source.
Its density is positive almost everywhere: the nonzero analytic
function $v_h$ on the physical line has isolated zeros; hence
$w_h>0$ almost everywhere, and its convolution powers are positive.
It follows that every nonzero polynomial has positive squared norm.
Thus all displayed moment and quotient forms are positive definite.
The original physical coordinate, both integration half-lines, and
the full root polynomial have remained in NCRP8.

Write $\mathbf F_{\mathcal I}$ for the increasing-power coefficient
matrix of $F_{\mathcal I}$. Its complete Gram and action are
\[
 \mathsf H_N=\mathbf F_{\mathcal I}^*G_N\mathbf F_{\mathcal I},
 \qquad
 \mathsf A=K_{\mathcal I}^{-1}CK_{\mathcal I}.
 \tag{NCRP9}
\]
The map $F_{\mathcal I}$ intertwines $\mathsf A$ and multiplication
by $S$ exactly, so
$\mathsf A w_\lambda=\lambda w_\lambda$.
Partition the blocks by $0,C,Z$, of sizes $v,q',m$. The actual
squared norm of the retained class is
\[
 \begin{split}
 \|v_\lambda\|_{G_N}^2={}&
 u_0^*\mathsf H_{00}u_0+\eta^*\mathsf H_{CC}\eta
       +z^*\mathsf H_{ZZ}z\\
 &+2\Re\{u_0^*\mathsf H_{0C}\eta
              +u_0^*\mathsf H_{0Z}z
              +\eta^*\mathsf H_{CZ}z\}.
 \end{split}                                                     \tag{NCRP10}
\]
These are the three diagonal terms and all six directed cross
pairings. There is no orthogonality claim about the DNE decomposition.

The full centered action has blocks
\[
 \mathsf D_{ij}
 =\sum_{r\in\{0,C,Z\}}\mathsf A_{ri}^*\mathsf H_{rj}
  +\sum_{r\in\{0,C,Z\}}\mathsf H_{ir}\mathsf A_{rj}
  -a\mathsf H_{ij}.
 \tag{NCRP11}
\]
Direct multiplication gives
$\mathsf D=\mathsf A^*\mathsf H+\mathsf H\mathsf A-a\mathsf H$.
For $\lambda=a\rho=c+a\delta+ia\gamma$, the original same-class
Hermitian action therefore obeys the exact full receiver
\[
 \begin{split}
 &u_0^*\mathsf D_{00}u_0+\eta^*\mathsf D_{CC}\eta
       +z^*\mathsf D_{ZZ}z\\
 &\qquad+
 2\Re\{u_0^*\mathsf D_{0C}\eta+
                  u_0^*\mathsf D_{0Z}z+
                  \eta^*\mathsf D_{CZ}z\}
       =2a\delta\, w_\lambda^*\mathsf H_Nw_\lambda>0.
 \end{split}                                                     \tag{NCRP12}
\]
The proof substitutes $\mathsf Aw_\lambda=\lambda w_\lambda$;
both terms contribute their full $\lambda,\bar\lambda$.
This is a statement on the complete retained class. It has not
assigned a sign to one of its summands.

\subsection{Orthogonal completion with every original cross block}

It is useful to obtain an exact minimum-distance formula without
changing the original minor. Put
$F_\perp^{\,0}=[F_0,F_C]$ and $u=(u_0,\eta)^{\mathsf T}$.
In the coefficient formulas below the maps are expressed in the
same increasing-power frame used for $G_N$. Define
\[
 \begin{gathered}
 \Gamma=L_a^*G_NL_a,\qquad
 H_{ZC}=L_a^*G_NF_\perp^{\,0},\qquad
 H_{CC}=(F_\perp^{\,0})^*G_NF_\perp^{\,0},\\
 P_Z=L_a\Gamma^{-1}L_a^*G_N,\qquad
 E_C=(I-P_Z)F_\perp^{\,0},\\
 \Sigma=H_{CC}-H_{ZC}^*\Gamma^{-1}H_{ZC},\qquad
 \widehat z=z+\Gamma^{-1}H_{ZC}u .
 \end{gathered}                                                   \tag{NCRP13}
\]
Completing the positive square proves that $\Sigma\succ0$:
for $u\ne0$, $E_Cu=0$ would place $F_\perp^{\,0}u$ in
$\operatorname{im}L_a$, contradicting the invertibility of
the full frame in NCRP3. Since $P_Z$ is the exact $G_N$-orthogonal
projector, $L_a^*G_NE_C=0$ and $E_C^*G_NE_C=\Sigma$.
Substitution now gives
\[
 \begin{gathered}
 v_\lambda=L_a\widehat z+E_Cu,\qquad
 \|v_\lambda\|_{G_N}^2
       =\widehat z^*\Gamma\widehat z+u^*\Sigma u,\\
 \operatorname{dist}_{G_N}(v_\lambda,\operatorname{im}L_a)^2
       =u^*\Sigma u>0.
 \end{gathered}                                                   \tag{NCRP14}
\]
This is a quantitative exact repair of the omitted-class issue.
The cross blocks in NCRP10 have contributed to both $\widehat z$
and $\Sigma$; they were not set to zero.

Write the action in the original algebraic frame, now ordered by
$(C,Z)$, as
\[
 A F_\perp^{\,0}=F_\perp^{\,0}A_{CC}+L_aA_{ZC},\qquad
 A L_a=F_\perp^{\,0}A_{CZ}+L_aA_{ZZ}.
 \tag{NCRP15}
\]
These matrices are the corresponding blocks of NCRP9 and therefore
are explicitly calculated from $K_{\mathcal I}^{-1}CK_{\mathcal I}$.
In the metric-orthogonal full frame $[L_a,E_C]$ they give
\[
 \begin{gathered}
 A_{\rm nat}=\Gamma^{-1}L_a^*G_NAL_a
              =A_{ZZ}+\Gamma^{-1}H_{ZC}A_{CZ},\\
 R_N=AL_a-L_aA_{\rm nat}=E_CA_{CZ},\\
 B_N=\Gamma^{-1}L_a^*G_NAE_C,\qquad
 D_N^{C}=\Sigma^{-1}E_C^*G_NAE_C,\\
 [A]_{[L_a,E_C]}
       =\begin{pmatrix}A_{\rm nat}&B_N\\ A_{CZ}&D_N^{C}\end{pmatrix}.
 \end{gathered}                                                   \tag{NCRP16}
\]
The $D_N^{C}$ here is a complement action block, not the native
Gamma metric named $D_N$ in earlier source modules.
For the lower-left block, subtracting $P_ZAL_a$ gives
$E_CA_{CZ}$, and multiplication by the inverse of $\Sigma$
proves its coefficient is exactly $A_{CZ}$.
Thus the full eigenclass has the two coupled equations
\[
 (A_{\rm nat}-\lambda I)\widehat z+B_Nu=0,\qquad
 A_{CZ}\widehat z+(D_N^{C}-\lambda I)u=0 .
 \tag{NCRP17}
\]
These equations require no inverse of $D_N^C-\lambda I$;
they hold on every original rank stratum.
The original FCI39 residual is exactly the matrix $R_N$ in
NCRP16. In particular
\[
 R_N^*G_NR_N=A_{CZ}^*\Sigma A_{CZ},\qquad
 \operatorname{rank}R_N=\operatorname{rank}A_{CZ}
                         =\operatorname{rank}R_{\mathcal T}>0 .
 \tag{NCRP18}
\]
The last equality follows by taking the quotient of the two full
coefficient frames: $F_\perp^{\,0}$ identifies its coefficient space
with $E/\operatorname{im}L_a$, while the omitted residue-dual
coordinates give another isomorphism of that same quotient.
Their transition is invertible, so both represent the same action
map into that quotient and have the same rank.

Define
\[
 \begin{split}
 D_Z&=A_{\rm nat}^*\Gamma+\Gamma A_{\rm nat}-a\Gamma,\\
 D_C&=(D_N^{C})^*\Sigma+\Sigma D_N^{C}-a\Sigma,\\
 X_N&=\Gamma B_N+A_{CZ}^*\Sigma.
 \end{split}                                                     \tag{NCRP19}
\]
The complete centered same-class equation becomes
\[
 \boxed{\quad
 \widehat z^*D_Z\widehat z+
    2\Re(\widehat z^*X_Nu)+u^*D_Cu
 =2a\delta(\widehat z^*\Gamma\widehat z+u^*\Sigma u).
 \quad}                                                         \tag{NCRP20}
\]
It follows by multiplying the two full blocks in NCRP16 against
the block diagonal metric $\operatorname{diag}(\Gamma,\Sigma)$.
Both leakage directions, $\Gamma B_N$ and $A_{CZ}^*\Sigma$,
remain separately in $X_N$. No vanishing or sign is inferred from
the positivity of $\Gamma$ or $\Sigma$.

The original native coefficient metric is still
$D_{a,N-v}^{(s)}$, at centre $a/2-4$ and order $s=1,a$.
Its exact comparison with the arithmetic section is the positive
matrix
\[
 (D_{a,N-v}^{(s)})^{-1/2}\Gamma
                         (D_{a,N-v}^{(s)})^{-1/2}.
 \tag{NCRP21}
\]
Equations NCRP13--20 do not assign this comparison an isometry or
a bounded condition number. The full conductor square in DNE and
XIM remains the actual morphism connecting those two presentations.

\subsection{The consecutive chart and the first physical metric}

The consecutive chart is used only through the proved chart
transition in the preceding section. In that chart
$\operatorname{im}L_a=\mathbb C[S]_{\le m-1}$.
At $N=q-1$ the remainder quotient has no relation to minimize,
so $G_{q-1}=H_{q-1}^{\rm ar}$ and $m<q$ implies that $Sf$
is still the literal polynomial for each $f$ in this section.
For two such polynomials, the physical identity
$S+\bar S=a$ gives
\[
 \langle f,Sg\rangle_{\rm ar}+
 \langle Sf,g\rangle_{\rm ar}
       =a\langle f,g\rangle_{\rm ar}.
 \tag{NCRP22}
\]
Projecting $Sf,Sg$ onto the same section changes neither of these
pairings. Therefore $D_Z=0$ at this cutoff in this chart.
The full receiver is consequently the exact evaluated identity
\[
 2\Re(\widehat z^*X_{q-1}u)+u^*D_Cu
       =2a\delta\,\|v_\lambda\|_{G_{q-1}}^2>0.
 \tag{NCRP23}
\]
It identifies the entire centered action of the retained class with
the complement and both cross terms. A centered compression by itself
does not bound these terms.

For its residual, let $p_n(y)$ be the monic real orthogonal
polynomials of the original arithmetic measure, with actual norms
$\omega_n=\int_{\mathbb R}|p_n(y)|^2dm_a^{\rm ar}(y)>0$.
Set $\phi_n(S)=p_n((S-c)/i)/\sqrt{\omega_n}$.
Multiplication by the real variable $y$ in this orthonormal frame is
Hermitian. Its coefficient from $\phi_{m-1}$ to $\phi_m$ is
$\sqrt{\omega_m/\omega_{m-1}}$, by comparing the leading coefficients.
For $j<m-1$, $y\phi_j$ has degree below $m$ and has no omitted
component. This proves the full formula
\[
 R_{q-1}
  =i\sqrt{\omega_m/\omega_{m-1}}\,
                       \phi_m e_{m-1}^*,\qquad
 \|R_{q-1}\|=\sqrt{\omega_m/\omega_{m-1}} .
 \tag{NCRP24}
\]
This is the original orthogonal-polynomial recurrence coefficient,
whose conventional formulation is DLMF18.2(iv)
\cite{human:dlmf18}. The proof above fixes its phase and scale.
For the two Gamma measures, and only for them,
$\omega_{n,s}=(2\pi)^{s/2}n!(s/2)_n$, giving
\[
 \|R_{q-1}^{(s)}\|=\sqrt{m(m-1+s/2)},\quad s=1,a,\qquad
 \frac{\|R_{q-1}^{(1)}\|}{a}\to16,\quad
 \frac{\|R_{q-1}^{(a)}\|}{a}\to\sqrt{264}.
 \tag{NCRP25}
\]
The complete arithmetic ratio in NCRP24 has not been replaced by
either Gamma ratio.

For the original arithmetic metric and nonzero
$f\in\mathbb C[S]_{\le m-1}$,
$(A-\lambda)f=(S-\lambda)f$ is literal. Thus
\[
 \begin{split}
 \|(A-\lambda)f\|_{\rm ar}^2
 &=(a\delta)^2\|f\|_{\rm ar}^2\\
 &\quad+\int_{\mathbb R}(y-a\gamma)^2
                      |f(c+iy)|^2dm_a^{\rm ar}(y)
 >(a\delta)^2\|f\|_{\rm ar}^2 .
 \end{split}                                                     \tag{NCRP26}
\]
The strict inequality follows because the integrand can vanish
almost everywhere only if the nonzero polynomial $f$ does.
This is a statement in the consecutive comparison chart at the
first cutoff; no estimate is transported to another chart by
dropping its conductor transition.

\subsection{The original observation projection on the completed class}

Let $\Lambda:E\to B_{\rm obs}$ be any one of the original packet
observation maps, and let $I_K$ be its original full kernel frame.
Its arithmetic orthogonal projector and complement are
\[
 P_K=I_K(I_K^*G_NI_K)^{-1}I_K^*G_N,\qquad P_B=I-P_K.
 \tag{NCRP27}
\]
For a zero-dimensional kernel the first map is zero. For the full
kernel it is the identity. In all cases
$P_K^2=P_K$, $P_K^*G_N=G_NP_K$, and
$P_Kv+P_Bv=v$.
The literal projected vector is
\[
 P_Kv_\lambda
  =P_KF_0u_0+P_KF_C\eta+P_KL_az.
 \tag{NCRP28}
\]
This equation, with the same original $P_K$, is the receiving
correction for any subsequent kernel response. It never replaces
$\ker\Lambda$ by $\operatorname{im}L_a$.

Put $x_K=P_Kv_\lambda$, $x_B=P_Bv_\lambda$ and
$D=A^*G_N+G_NA-aG_N$. Exact substitution of
$Av_\lambda=\lambda v_\lambda$ gives
\[
 \begin{split}
 x_K^*Dx_K
   &=2a\delta\|x_K\|_{G_N}^2
                    -2\Re(x_K^*G_NAx_B),\\
 x_B^*Dx_B
   &=2a\delta\|x_B\|_{G_N}^2
                    -2\Re(x_B^*G_NAx_K),\\
 2\Re(x_K^*Dx_B)
   &=2\Re\{x_K^*G_NAx_B+x_B^*G_NAx_K\}.
 \end{split}                                                     \tag{NCRP29}
\]
For example, $Ax_K=\lambda(x_K+x_B)-Ax_B$ and
$x_K^*G_Nx_B=0$ prove the first line. The second follows by
the other decomposition; in the last line expand $D$ and use
orthogonality to remove only its actual $aG_N$ term.
Adding all three equations returns
$2a\delta\|v_\lambda\|_{G_N}^2$ exactly.
Thus both projected signs still depend on their actual leakage
pairings, now evaluated on the completed class NCRP28.

\subsection{The literal extension of the relative comparison cone}

For four copies let $\mathcal H=\mathbb C^{4m}$ and let
$L=I_4\otimes L_a$. The original coefficient cone is
$B\to W\oplus S\to\mathcal H$ with differentials
$b\mapsto(b,-b)$ and $(w,s)\mapsto w+s$.
The enlarged arithmetic packet cone is obtained by replacing its
last term with $E^{\oplus4}$ and using $L(w+s)$ there.
Since $L$ is injective, its middle cohomology remains
$(S\cap W)/B$. If $T=S+W$, its final cohomology is
$E^{\oplus4}/L(T)$, and the exact sequence is
\[
 0\longrightarrow L(\mathcal H)/L(T)
 \longrightarrow E^{\oplus4}/L(T)
 \longrightarrow E^{\oplus4}/L(\mathcal H)
 \longrightarrow0.
 \tag{NCRP30}
\]
Each map is induced by the identity on representatives. Its exactness
follows by testing whether a representative lies in $L(\mathcal H)$.
The full complement from NCRP3 provides explicit representatives for
the last quotient, of dimension $4(v+q')$.
All quotient metrics are the attained quotients of the full original
$G_N^{\oplus4}$, so their Schur blocks retain the cross terms in
NCRP10 and NCRP13. A retained eigenclass in one indicated copy lies
outside $L(\mathcal H)$ and therefore has a nonzero class in this
enlarged final cohomology.

Multiplication by $S$ on the full packet is the exact action
NCRP9/NCRP16. The native subspace has the nonzero defect NCRP18,
so the action is carried by the complete ambient block matrices.
No unproved preservation of $B,S,W,T$ is used to create an action
on their quotients. In particular NCRP30 is an explicit extension
of the old relative cohomology, with its additional support and
metric; it does not identify a relative denominator with a theta
boundary. The original native determinant allocations remain valid
on their own maps and metrics, and their arithmetic use proceeds
through the completed class and the complete receiving equations above.


% Additive proof fragment.  All matrices use literal powers of S.
% Provider bibliography key: prog:arithmetic.
\section{The full arithmetic action boundary in the original quotient}

The polynomial quotient and positive arithmetic moments in this section are
the original ones of the arithmetic construction
\cite{prog:arithmetic}.  The calculation below uses the full relation ideal,
the original physical line, and the attained quotient metric.  Every identity
needed for the boundary calculation is proved here.

\subsection{Spaces, moments, and the attained section}

In the actual packet \(a=k+4\), with \(k\) in its original guarded domain;
retain that domain and its original \(0<\delta<1/2\) and \(\gamma>2\).
The finite identities below in fact use only that \(a\) is a positive
integer and \(\delta,\gamma>0\).  Put
\begin{equation}
 \beta_{rt}=\frac a2+(2r-a)\delta+i(2t-a)\gamma,\qquad
 \chi(S)=\prod_{r=0}^{a}\prod_{t=0}^{a}(S-\beta_{rt}),\qquad
 q=(a+1)^2 .
 \tag{NAB1}\label{eq:NAB1}
\end{equation}
These roots are pairwise distinct: equality of real parts gives equality
of \(r\), and equality of imaginary parts then gives equality of \(t\).
No root or coefficient of \(\chi\) is discarded.  Write
\(\mathcal P_j=\{p\in\mathbb C[S]:\deg p\leq j\}\), with
\(\mathcal P_{-1}=\{0\}\), and
\(\mathcal E=\mathbb C[S]/(\chi)\).  Coordinates in
\(\mathcal P_j\) are the literal coefficients of
\(1,S,\ldots,S^j\); coordinates in \(\mathcal E\) are the corresponding
classes of \(1,S,\ldots,S^{q-1}\).

Fix an integer \(N\geq q-1\).  Use the original positive measure
\(d\mu_N(y)\) on the physical line \(S=a/2+iy\), with all of its original
components and both integration regions.  The subscript allows the measure
to depend on the fixed endpoint \(N\).  Within this calculation the measure
is held fixed.  Its moment matrix through degree \(N+1\) is
\begin{equation}
 H_j^{[N]}=
 \left(
 \int_{\mathbb R}
    \overline{(a/2+iy)^r}(a/2+iy)^t\,d\mu_N(y)
 \right)_{0\leq r,t\leq j},
 \qquad 0\leq j\leq N+1 .
 \tag{NAB2}\label{eq:NAB2}
\end{equation}
All moments displayed here are finite and \(H_{N+1}^{[N]}\) is positive
definite in the original arithmetic construction.  Equivalently, the
integral of \(|p(a/2+iy)|^2\) is strictly positive for every nonzero
\(p\in\mathcal P_{N+1}\).  The superscript records that
\(H_N^{[N]}\) and \(H_{N+1}^{[N]}\) use the same measure; an endpoint change
of measure is not implicit in this notation.  In the remainder of this
section write these matrices as \(H_N\) and \(H_{N+1}\).

The Hermitian product is conjugate linear in its first entry:
\(\langle f,g\rangle_N=\int\overline{f(a/2+iy)}g(a/2+iy)d\mu_N(y)\).
For \(f,g\in\mathcal P_N\), the identity
\(\overline S+S=a\) on this actual line gives
\begin{equation}
 \langle Sf,g\rangle_N+\langle f,Sg\rangle_N
       =a\langle f,g\rangle_N .
 \tag{NAB3}\label{eq:NAB3}
\end{equation}
The two left-hand products use the moment matrix through degree \(N+1\).
Thus (NAB3) retains the physical coordinate rather than replacing the
action by an abstract self-adjoint multiplication operator.

Let \(\rho_N:\mathcal P_N\longrightarrow\mathcal E\) be reduction modulo
\(\chi\).  Polynomial division by the original monic \(\chi\) proves
\begin{equation}
 \ker\rho_N=\mathcal R_N:=\chi\mathcal P_{N-q},\qquad
 \mathcal P_N/\mathcal R_N \xrightarrow[\rho_N]{\ \cong\ }\mathcal E .
 \tag{NAB4}\label{eq:NAB4}
\end{equation}
For \(N=q-1\), \(\mathcal R_N=\{0\}\).
Let \(E_N:\mathbb C^q\to\mathbb C^{N+1}\) be the coefficient inclusion of
\(\mathcal P_{q-1}\) into \(\mathcal P_N\).  If \(N\geq q\), let \(W_N\)
have columns the coefficient vectors of
\(\chi,S\chi,\ldots,S^{N-q}\chi\).  Its columns are linearly independent
because multiplication by a nonzero polynomial is injective.
Consequently \(W_N^*H_NW_N\) is positive definite.  The attained section is
\begin{equation}
 T_N=E_N-W_N(W_N^*H_NW_N)^{-1}W_N^*H_NE_N,\qquad
 G_N=T_N^*H_NT_N ;
 \quad T_{q-1}=E_{q-1}.
 \tag{NAB5}\label{eq:NAB5}
\end{equation}
Here \(T_N:\mathcal E\to\mathcal P_N\) uses the fixed coefficient
coordinates, and \(G_N\) is its \(q\)-by-\(q\) attained metric.
Indeed, \(\rho_NT_N=I_{\mathcal E}\) by (NAB4), and direct multiplication
gives \(W_N^*H_NT_N=0\).  Every representative of \(v\in\mathcal E\)
is \(T_Nv+W_Nz\) for a unique \(z\).  Orthogonality therefore proves
\[
 \|T_Nv+W_Nz\|_N^2
 =v^*G_Nv+z^*W_N^*H_NW_Nz.
\]
This proves attainment, uniqueness, and the original minimum property
without a change of relation space.  It also proves \(G_N>0\), since
\(T_Nv\ne0\) whenever \(v\ne0\).

\subsection{The exact one-column boundary and its Hermitian form}

Let \(A:\mathcal E\to\mathcal E\) be multiplication by the original \(S\)
modulo \(\chi\).  Write
\(\chi(S)=S^q+\sum_{j=0}^{q-1}c_jS^j\).
In the fixed ordered basis,
\begin{equation}
 A e_j=e_{j+1}\quad(0\leq j<q-1),\qquad
 A e_{q-1}=-\sum_{j=0}^{q-1}c_je_j .
 \tag{NAB6}\label{eq:NAB6}
\end{equation}
Embed \(T_N\) into \(\mathcal P_{N+1}\) by appending a zero top
coefficient, writing this embedding as \(\widetilde T_N\).
Define the row, polynomial column, and vector
\begin{equation}
 \ell_N=[S^N]T_N:\mathcal E\longrightarrow\mathbb C,\qquad
 b_N=\operatorname{coeff}_{0,\ldots,N+1}
             \bigl(\chi(S)S^{N+1-q}\bigr),\qquad
 \xi_N=\widetilde T_N^*H_{N+1}b_N\in\mathbb C^q .
 \tag{NAB7}\label{eq:NAB7}
\end{equation}
The notation \([S^N]\) means the literal coefficient functional, with no
factorial, root evaluation, or rescaling.

Let \(M_N:\mathcal P_N\to\mathcal P_{N+1}\) denote multiplication by \(S\)
in those same coordinates.  There is a linear map
\(Z_N:\mathcal E\to\mathcal R_N\), embedded in \(\mathcal P_{N+1}\),
such that
\begin{equation}
 M_NT_N-\widetilde T_N A=b_N\ell_N+Z_N,\qquad
 \widetilde T_N^*H_{N+1}Z_N=0.
 \tag{NAB8}\label{eq:NAB8}
\end{equation}
To prove the first identity on \(v\in\mathcal E\), reduce
\(S(T_Nv)-T_N(Av)\) modulo \(\chi\).
The result is \(Av-Av=0\), so the polynomial is divisible by \(\chi\).
Its degree is at most \(N+1\), and its coefficient of \(S^{N+1}\) is
exactly \(\ell_Nv\).  Monicity of the full \(\chi\) now shows that
subtracting \(\chi S^{N+1-q}\ell_Nv\) leaves an element of
\(\chi\mathcal P_{N-q}\).
This subtraction is linear in \(v\), giving \(Z_N\).
If \(N=q-1\), the remainder has degree less than \(q\) and is divisible
by \(\chi\), hence is zero.
Otherwise the second identity is exactly the orthogonality in (NAB5):
both factors of this product have degree at most \(N\), so the use of
\(H_{N+1}\) restricts to the same \(H_N\).

Define the full arithmetic action form
\begin{equation}
 D_N=A^*G_N+G_NA-aG_N .
 \tag{NAB9}\label{eq:NAB9}
\end{equation}
It has the exact boundary formula
\begin{equation}
 \boxed{\displaystyle
 D_N=-\xi_N\ell_N-\ell_N^*\xi_N^*.}
 \tag{NAB10}\label{eq:NAB10}
\end{equation}
For a direct proof, (NAB8) and its orthogonality give
\[
 \widetilde T_N^*H_{N+1}M_NT_N
       =G_NA+\xi_N\ell_N .
\]
Taking conjugate transposes gives the other term.
Applying (NAB3) to every pair of columns of \(T_N\) gives
\[
 aG_N=
 \widetilde T_N^*H_{N+1}M_NT_N+
 (M_NT_N)^*H_{N+1}\widetilde T_N
 =G_NA+A^*G_N+\xi_N\ell_N+\ell_N^*\xi_N^*.
\]
Rearranging proves (NAB10), including both minus signs.
In particular \(\operatorname{rank}D_N\leq2\).
No estimate for a trial section or different denominator was used.

\subsection{Exact inertia, trace, and boundary magnitude}

For every root \(\beta\) in (NAB1), define the original Lagrange
polynomial and its class
\begin{equation}
 L_\beta(S)=\frac{\chi(S)}{(S-\beta)\chi'(\beta)}
       \in\mathcal P_{q-1},\qquad
 v_\beta=[L_\beta]\in\mathcal E .
 \tag{NAB11}\label{eq:NAB11}
\end{equation}
The denominator \(\chi'(\beta)\) is nonzero by the distinctness proved
after (NAB1).  At the original roots \(L_\beta(\beta)=1\) and
\(L_\beta(\beta')=0\) for \(\beta'\ne\beta\).
Thus \(v_\beta\ne0\), and the \(q\) such classes form a basis: a linear
relation among them evaluates at each root to its corresponding
coefficient.  The literal polynomial identity
\((S-\beta)L_\beta=\chi/\chi'(\beta)\) proves the exact morphism
\begin{equation}
 Av_\beta=\beta v_\beta,\qquad
 v_{\beta_{rt}}^*D_Nv_{\beta_{rt}}
      =2(2r-a)\delta\,v_{\beta_{rt}}^*G_Nv_{\beta_{rt}} .
 \tag{NAB12}\label{eq:NAB12}
\end{equation}
In particular the expression is strictly negative at \(r=0\) and
strictly positive at \(r=a\).
A Hermitian form of rank at most two with both a positive and a negative
value must have rank exactly two and one sign of each kind:
\begin{equation}
 \operatorname{inertia}(D_N)=(1,1,q-2).
 \tag{NAB13}\label{eq:NAB13}
\end{equation}
For completeness, this elementary inference can be seen by diagonalizing
the Hermitian form using successive orthogonal eigenvectors.  Such a
basis is obtained by maximizing its real Rayleigh quotient on the compact
unit sphere; differentiation in every real and imaginary tangent
direction makes a maximizing vector an eigenvector.  Its orthogonal
complement is invariant because the matrix is Hermitian, so induction
gives the diagonalization.  A positive value requires a positive
diagonal entry and a negative value requires a negative diagonal entry.
The rank bound allows precisely those two nonzero entries.

The full original grid also gives an exact trace cancellation:
\begin{equation}
 \operatorname{tr}(G_N^{-1}D_N)
 =\overline{\operatorname{tr}A}+\operatorname{tr}A-aq
 =2\operatorname{Re}\sum_{r,t=0}^{a}\beta_{rt}-aq=0 .
 \tag{NAB14}\label{eq:NAB14}
\end{equation}
The first equality uses only cyclicity of finite matrix trace.
For the last equality,
\(\sum_{r=0}^a(2r-a)=0=\sum_{t=0}^a(2t-a)\), so the root sum
is exactly \(aq/2\).
The Lagrange basis in (NAB11) proves that this root sum is the trace of
the actual \(A\); no approximate eigenvalue calculation is involved.

Take the positive Hermitian square root of \(G_N\), and define
\begin{equation}
 B_N=G_N^{-1/2}D_NG_N^{-1/2},\qquad
 u_N=G_N^{-1/2}\xi_N,\qquad
 w_N=G_N^{-1/2}\ell_N^*,\qquad
 c_N=w_N^*u_N=\ell_NG_N^{-1}\xi_N .
 \tag{NAB15}\label{eq:NAB15}
\end{equation}
The square root is obtained by the same finite Hermitian
diagonalization, taking positive square roots of the positive entries of
\(G_N\).  Equations (NAB10) and (NAB14) prove
\begin{equation}
 B_N=-u_Nw_N^*-w_Nu_N^*,\qquad
 -2\operatorname{Re}c_N=\operatorname{tr}B_N=0.
 \tag{NAB16}\label{eq:NAB16}
\end{equation}
Congruence by the invertible \(G_N^{-1/2}\) preserves positive and
negative subspaces, hence (NAB13) applies to \(B_N\) too.
In particular \(u_N,w_N\) are linearly independent.
Set \(b=\|w_N\|>0\), \(e_1=w_N/b\), and write
\[
 u_N=(c_N/b)e_1+\eta e_2,\qquad
 \eta=\left(\|u_N\|^2-|c_N|^2/b^2\right)^{1/2}>0,
\]
where \(e_2\) is the resulting unit vector orthogonal to \(e_1\).
On their span the matrix of \(B_N\) is exactly
\[
 \begin{pmatrix}-2\operatorname{Re}c_N&-b\eta\\
                 -b\eta&0\end{pmatrix},
\]
and on its orthogonal complement it is zero.  Since
\(\operatorname{Re}c_N=0\), its nonzero eigenvalues are
\begin{equation}
 \boxed{\displaystyle
 \lambda_\pm(B_N)=\pm\kappa_N,\qquad
 \kappa_N=
 \sqrt{\|u_N\|^2\|w_N\|^2-
                    (\operatorname{Im}(w_N^*u_N))^2}>0 .}
 \tag{NAB17}\label{eq:NAB17}
\end{equation}
The equivalent formula using only the original coefficient matrices is
\begin{equation}
 \kappa_N^2=
  (\xi_N^*G_N^{-1}\xi_N)
  (\ell_NG_N^{-1}\ell_N^*)
  -|\ell_NG_N^{-1}\xi_N|^2
  =\tfrac12\operatorname{tr}\bigl((G_N^{-1}D_N)^2\bigr).
 \tag{NAB18}\label{eq:NAB18}
\end{equation}
The first equality follows from (NAB15)--(NAB17), and the second follows
because \(G_N^{-1}D_N\) is similar to \(B_N\).
The exact grid eigenclasses in (NAB12) furthermore give the finite bound
\begin{equation}
 \kappa_N\geq2a\delta .
 \tag{NAB19}\label{eq:NAB19}
\end{equation}
Indeed, put \(z=G_N^{1/2}v_{\beta_{0t}}\).  It is nonzero and
\(z^*B_Nz/(z^*z)=-2a\delta\).
Every Rayleigh quotient of the diagonal matrix with entries
\(-\kappa_N,\kappa_N,0,\ldots,0\) lies in
\([-\kappa_N,\kappa_N]\); this proves (NAB19).
The choice \(r=a\) gives the matching positive quotient \(2a\delta\).

If the complete grid trace cancellation is not available, the boundary
identity (NAB10) still has its stated sign whenever (NAB2)--(NAB8) hold,
but the two eigenvalues are instead
\[
 -\operatorname{Re}c_N
 \ \pm\sqrt{\|u_N\|^2\|w_N\|^2-(\operatorname{Im}c_N)^2}.
\]
This follows by taking the characteristic polynomial of the displayed
two-by-two matrix.  In the actual full grid the shift vanishes exactly
by (NAB14).  Thus omission of a root subfamily cannot silently retain
the full-grid trace cancellation.

\subsection{An exact original-coordinate sign diagnostic}

The following small example is an exact algebraic diagnostic only.  It
is not an actual packet or period, and is not an evaluation of the
arithmetic source.  It neither replaces nor relaxes any original
large-degree finite guard.
Take \(a=1\), \(\delta=1/4\), \(\gamma=3\), \(q=4\), \(N=4\),
\(S=1/2+iy\), and the standard real Gaussian measure.  The original
root polynomial and relation space are
\begin{equation}
 \begin{aligned}
 \chi(S)&=
 \bigl((S-\tfrac12-\tfrac14)^2+9\bigr)
 \bigl((S-\tfrac12+\tfrac14)^2+9\bigr)\\
 &=S^4-2S^3+\tfrac{155}{8}S^2-\tfrac{147}{8}S+\tfrac{22185}{256},
 \qquad \mathcal R_4=\operatorname{span}\{\chi\}.
 \end{aligned}
 \tag{NAB20}\label{eq:NAB20}
\end{equation}
The roots are exactly \(1/2\pm1/4\pm3i\), all in the right half-plane.
For \(m\geq0\), put
\(\mu_{2m}=(2m-1)!!\), with \((-1)!!=1\), and
\(\mu_{2m+1}=0\).  These are the Gaussian moments: the odd integral
vanishes under \(y\mapsto-y\), and integration by parts gives
\(\mu_{2m}=(2m-1)\mu_{2m-2}\).
The moment entries, entirely in the original \(S\)-coordinates, are
\begin{equation}
 (H_5)_{mn}
   =\sum_{r=0}^{m}\sum_{t=0}^{n}
    \binom mr\binom nt\,2^{-(m+n-r-t)}
            (-\mathrm i)^r\mathrm i^t\mu_{r+t}.
 \tag{NAB21}\label{eq:NAB21}
\end{equation}
Here \(\mathrm i^2=-1\), and \(0\leq m,n\leq5\).
Every nonzero polynomial has positive squared integral because a
nonzero polynomial has finitely many zeros and the Gaussian density is
strictly positive.  Thus the moment and quotient matrices are positive
definite.

Write
\(c=(22185/256,-147/8,155/8,-2,1)^{\mathsf T}\), and put
\(d=316481153\).
The finite sums (NAB21) give \(c^*H_4c=d/65536\).
Substitution into (NAB5) and (NAB7) gives, exactly,
\begin{equation}
 \begin{aligned}
 \ell_4&=\frac1d(-4407552,-2203776,1749184,3725664),\\
 \xi_4&=\frac1{1024}(0,44548,44548,-51849)^{\mathsf T}.
 \end{aligned}
 \tag{NAB22}\label{eq:NAB22}
\end{equation}
and hence
\begin{equation}
 D_4=-\xi_4\ell_4-\ell_4^{\mathsf T}\xi_4^{\mathsf T}.
 \tag{NAB23}\label{eq:NAB23}
\end{equation}
To specify the remaining finite computation without omitting its
inputs, take \(E_4\) to be the first four columns of the five-by-five
identity, \(T_4=E_4-(65536/d)c(c^*H_4E_4)\), and
\(G_4=T_4^*H_4T_4\), with \(H_4\) the leading block of (NAB21).
Direct substitution and ordinary exact matrix multiplication give
\begin{equation}
 \begin{aligned}
 \operatorname{tr}(G_4^{-1}D_4)&=0,\\
 \frac{\det(zG_4-D_4)}{\det G_4}
      &=z^2\left(z^2-\frac{46738154483830145}{497781812232192}\right),\\
 \kappa_4^2&=\frac{46738154483830145}{497781812232192}.
 \end{aligned}
 \tag{NAB24}\label{eq:NAB24}
\end{equation}
All the matrices and exact rational substitutions for this diagnostic
are recorded by the accompanying
\texttt{exact\_boundary\_example.py} and
\texttt{EXACT\_EXAMPLE.json}.  The general proof is (NAB3)--(NAB19);
the example supplies an independent check of its signs and its use of
the nontrivial attained relation space.

\paragraph{Scope of the calculation.}
Equations (NAB10), (NAB13), and (NAB17)--(NAB19) determine the full
action form and its full-space inertia.  For a specified map
\(J:\mathbb C^d\to\mathcal E\), their exact receiver is
\[
 J^*D_NJ=-(J^*\xi_N)(\ell_NJ)
                 -(\ell_NJ)^*(\xi_N^*J).
\]
This formula carries the original class-position data into the
restricted form.  Evaluating the sign on an individual projected class
requires its actual \(J\), or the corresponding attained projection,
inside that formula.  No individual projected sign, absolute
allocation, or value of the responses \(U,V\) is asserted here.


% Exact FCI/DNE-to-EC dictionary; new actual class completion.
\section{The original EC three-column response after full-class completion}

We use exactly the arithmetic packet, physical source and attained
section of NCRP8. The letter $a$ is the tensor degree used in this
manuscript; it replaces EC's tensor-degree letter $k$, and does not
change the original source order or packet. Let $p_n(S)$ be the
original arithmetic monic orthogonal polynomials at $S=c+iy$ and
let $\omega_n=\|p_n\|_{\rm ar}^2$. Explicitly, if
$p_n^{y}$ is the monic real-variable polynomial in NCRP24, then
$p_n(S)=i^n p_n^{y}((S-c)/i)$; this is monic in the original $S$
and has exactly the same norm $\omega_n$. Put $b_n=[p_n]\in E$.
Here $b_n$ retains EC's polynomial-class notation; the relation
column called $b_N$ in NAB7 is called $r_N$ below.
At each original cutoff $N$ the three columns and full Grams are
precisely
\[
 \begin{gathered}
 W=[b_N,b_{N+1},v_\lambda],\qquad F=W^*G_NW,\\
 K=W^*G_N I_K(I_K^*G_NI_K)^{-1}I_K^*G_NW,\qquad B=F-K,\\
 F_{13}=b_N^*G_Nv_\lambda,\quad
 F_{23}=b_{N+1}^*G_Nv_\lambda,\quad F_{33}=\|v_\lambda\|_{G_N}^2 .
 \end{gathered}                                                   \tag{NCE1}
\]
These are EC8 and EC37 \cite{prog:EC}; the kernel projector is
the original $P_K$ of NCRP27, with its complete inverse.

\subsection{Identify both boundary rows, with their orientation}

Let $\ell_N,\xi_N$ be the full action-boundary rows proved in the
preceding section, and put $r_N(S)=\chi(S)S^{N+1-q}$.
For every class $x$, its attained lift $T_Nx$ has leading coefficient
$\ell_Nx$. Because $p_N$ is monic and orthogonal to every lower
degree polynomial,
\[
 \langle p_N,T_Nx\rangle_{\rm ar}=\omega_N\ell_Nx.
 \tag{NCE2}
\]
The polynomial $p_N-T_Nb_N$ is an old relation, orthogonal to
$T_Nx$. Thus the left side is $b_N^*G_Nx$, and consequently
\[
 \boxed{\ell_N=\omega_N^{-1}b_N^*G_N.}
 \tag{NCE3}
\]
Since $p_{N+1}$ and $r_N$ are both monic of degree $N+1$,
$p_{N+1}-r_N$ has degree at most $N$ and represents $b_{N+1}$.
The orthogonality of $p_{N+1}$ to every polynomial of degree
at most $N$ gives
\[
 \langle T_Nx,r_N\rangle_{\rm ar}
   =-\langle T_Nx,p_{N+1}-r_N\rangle_{\rm ar}
   =-x^*G_Nb_{N+1}.
 \tag{NCE4}
\]
The definition of $\xi_N$ therefore proves the second exact row:
\[
 \boxed{\xi_N=-G_Nb_{N+1}.}
 \tag{NCE5}
\]
In particular the original centered action matrix is
\[
 \boxed{\quad
 D_N=A^*G_N+G_NA-aG_N
 =\frac{G_Nb_{N+1}b_N^*G_N+G_Nb_Nb_{N+1}^*G_N}{\omega_N}.
 \quad}                                                         \tag{NCE6}
\]
This is a direct coefficient proof of the rank-two marked-pairing
matrix on the original arithmetic packet. Its sign agrees with EC10.

To identify the complex coefficient, put $v=v_\lambda$,
$f=T_Nv$, $E=F_{33}$ and
\[
 \mu_{N,\lambda}
   =\frac{\int_{\mathbb R}y|f(c+iy)|^2dm_a^{\rm ar}(y)}{E}.
 \tag{NCE7}
\]
The polynomial $(S-\lambda)f$ is a relation through degree $N+1$.
Its top coefficient is $\ell_Nv$, and subtracting
$(\ell_Nv)r_N$ leaves an old relation. Pairing with the attained
section eliminates exactly that old relation and gives
\[
 (v^*\xi_N)(\ell_Nv)
   =\langle f,(S-\lambda)f\rangle_{\rm ar}
   =E\{-a\delta+i(\mu_{N,\lambda}-a\gamma)\}.
 \tag{NCE8}
\]
Combining NCE3 and NCE5, the left side is
$-\overline{F_{23}}F_{13}/\omega_N$. Taking its conjugate proves
the literal EC current, including its phase:
\[
 \boxed{\quad
 \mathfrak c_{N,\lambda}
 :=\frac{\overline{F_{13}}F_{23}}{\omega_NE}
 =a\delta+i(\mu_{N,\lambda}-a\gamma).
 \quad}                                                         \tag{NCE9}
\]
Its positive real part proves $F_{13}\ne0$ and $F_{23}\ne0$.
No asymptotic estimate or period minor is needed to justify these
two denominators.

\subsection{The exact finite current bound in the original source}

Here is a direct finite proof of the useful strict current interval.
The arithmetic source is even in $y$: conjugation preserves the
complete quartet and $2\xi/h$, so $w_h(-y)=w_h(y)$, and its
convolution powers remain even. The full physical root polynomial
$i^{-q}\chi(c+iy)$ is real and has parity $(-1)^q$, because its
roots in the $y$-coordinate are invariant under conjugation and negation.
In particular $\chi(c-iy)=(-1)^q\chi(c+iy)$; no assertion that
$q$ is even is used.
Hence reflection preserves the relation space and its orthogonal
complement $\mathcal V_N=\operatorname{im}T_N$.

Let $\Pi_N$ be the source-orthogonal projector onto $\mathcal V_N$,
let $J=\Pi_NM_y|_{\mathcal V_N}$, and retain
$\widehat b_n=T_Nb_n$. The operator $J$ is Hermitian and
reflection anticommutes with it. The two vectors
$\widehat b_N,\widehat b_{N+1}$ have opposite reflection parities,
so their actual inner product is zero.
Multiplying an attained representative by $S$, subtracting its
leading multiple of $p_{N+1}$, and projecting gives the full
action identity
\[
 T_NAT_N^{-1}
   =cI+iJ+
       \frac{\widehat b_{N+1}\widehat b_N^*}{\omega_N}.
 \tag{NCE10}
\]
The proof of the last coefficient is NCE2; the class of the
subtracted polynomial is $b_{N+1}$, which is added back after
projection. Thus the term has the displayed positive sign.
The old relations are orthogonal to $\mathcal V_N$ and no
additional term is removed.

Put $d=a\delta>0$, $t=a\gamma>0$ and
$R=(d+it-iJ)^{-1}$. For every vector $x$,
$\|(d+it-iJ)x\|\,\|x\|\ge
\Re\langle x,(d+it-iJ)x\rangle=d\|x\|^2$.
The finite square matrix is injective, hence invertible, with
$\|R\|\le d^{-1}$.
The eigenclass equation and NCE10 imply
\[
 f=\frac{F_{13}}{\omega_N}R\widehat b_{N+1},\qquad
 \mathfrak c_{N,\lambda}
  =\frac{\langle\widehat b_{N+1},R\widehat b_{N+1}\rangle}
         {\|R\widehat b_{N+1}\|^2}.
 \tag{NCE11}
\]
The second equation follows by substituting the first into
$F_{23}$ and $E=\|f\|^2$. The original vector and source mass
have not been rescaled.

The spectral measure of the Hermitian $J$ at
$\widehat b_{N+1}$ is even, since reflection anticommutes with
$J$ and preserves that vector up to sign. Under the resolvent
squared norm, the paired spectral points $(x,-x)$ have mean
\[
 \frac{x\{[d^2+(t-x)^2]^{-1}-[d^2+(t+x)^2]^{-1}\}}
 {[d^2+(t-x)^2]^{-1}+[d^2+(t+x)^2]^{-1}}
       =\frac{2tx^2}{d^2+t^2+x^2}\in[0,2t).
 \tag{NCE12}
\]
There is some nonzero spectral mass away from zero. Otherwise
$R\widehat b_{N+1}$ would be a scalar multiple of
$\widehat b_{N+1}$, orthogonal to $\widehat b_N$; taking the
$\widehat b_N$ inner product in NCE11 would contradict
$F_{13}\ne0$.
Thus its total weighted mean is strictly between zero and $2t$.
This total mean is $\mu_{N,\lambda}$ by NCE11 and the definition
of $J$. Consequently
\[
 0<\mu_{N,\lambda}<2a\gamma,\qquad
 \Re\mathfrak c_{N,\lambda}=a\delta,\qquad
 |\Im\mathfrak c_{N,\lambda}|<a\gamma,\qquad
 |\mathfrak c_{N,\lambda}|<a\sqrt{\delta^2+\gamma^2}.
 \tag{NCE13}
\]
The finite spectral argument is the one behind EC11--13. It is
reproduced here to keep the physical moment and orientation explicit.
It supplies no estimate for the actual observation responses.

\subsection{The same original ratios and all three projected terms}

By NCE3--5 the exact dictionary, with no row exchanged, is
\[
 \boxed{\quad
 U=\frac{K_{13}}{F_{13}}
     =\frac{\ell_NP_Kv}{\ell_Nv},\qquad
 V=\frac{K_{23}}{F_{23}}
     =\frac{\xi_N^*P_Kv}{\xi_N^*v}.
 \quad}                                                         \tag{NCE14}
\]
Both occurrences of the minus sign in NCE5 cancel in the second
ratio. These are precisely EC37's $U,V$.
Set
$P_K^{\rm resp}=\bar UV$,
$P_B^{\rm resp}=(1-\bar U)(1-V)$ and
$P_\times^{\rm resp}=\bar U+V-2\bar UV$.
For $x_K=P_Kv$ and $x_B=(I-P_K)v$, define the original three
pairings by
$\Phi_K=x_K^*D_Nx_K$,
$\Phi_B=x_B^*D_Nx_B$ and
$\Phi_\times=2\Re(x_K^*D_Nx_B)$.
Substitution in NCE6 gives
\[
 \boxed{\frac{\Phi_X}{2E}
       =\Re(\mathfrak c_{N,\lambda}P_X^{\rm resp}),
           \qquad X=K,B,\times.}
 \tag{NCE15}
\]
For example the kernel term is
$2\Re(\overline{K_{13}}K_{23})/\omega_N$.
The mixed term is
$2\Re(\overline{K_{13}}B_{23}+\overline{B_{13}}K_{23})/\omega_N$;
substitution of the two ratios gives exactly the displayed
$P_\times^{\rm resp}$. Their sum is one, so the sum of the three
pairings is exactly $2a\delta E$.

Let $w_\lambda=(u_0,\eta,z)$ be NCRP5 and let the two exact rows be
$r_1=\ell_N$, $r_2=\xi_N^*$.
For $j=1,2$, the numerator of the corresponding response is
\[
 r_jP_Kv
    =r_jP_KF_0u_0+r_jP_KF_C\eta+r_jP_KL_az,
 \qquad
 r_jv=r_jF_0u_0+r_jF_C\eta+r_jL_az .
 \tag{NCE16}
\]
Thus NCE14 evaluates the full coefficient map needed to pass from
the native construction to the actual EC responses. Dropping the
first two terms in either expression changes the class and the
response. The individual products and their imaginary parts are
still those of these complete ratios.

Finally let $\theta_K=\|P_Kv\|_{G_N}^2/E$. The two separate action
leakages obtained in NCRP29 are
\[
 \begin{split}
 \frac{2\Re(x_K^*G_NAx_B)}{2E}
       &=a\delta\theta_K-\Re(\mathfrak c P_K^{\rm resp}),\\
 \frac{2\Re(x_B^*G_NAx_K)}{2E}
       &=a\delta(1-\theta_K)-\Re(\mathfrak c P_B^{\rm resp}).
 \end{split}                                                     \tag{NCE17}
\]
These are exactly EC41 with its two block directions retained.
Their sum is $\Phi_\times/(2E)$.
The more precise parity-dependent EC current, and its original
eventual guards and analytic error radius, retain their existing
verification status. This source proves the actual row dictionary,
the complete-class input, the finite current interval, and the
exact three-term and two-leakage receiving maps; it does not
reassign their period-dependent response products or re-prove
the later EC asymptotic radius.

\subsection{The exact parity reduction of the three-column Gram}

The same original reflection used in NCE10 proves a further evaluated
entry of EC's actual Gram. Since the two attained columns have opposite
parities, reflection is a unitary map and changes the sign of their
inner product. Thus
\[
 F_{12}=0,\qquad
 \ell_NG_N^{-1}\xi_N=-F_{12}/\omega_N=0,\qquad
 \kappa_N=\frac{\sqrt{F_{11}F_{22}}}{\omega_N}.
 \tag{NCE18}
\]
The last equality follows by substituting NCE3 and NCE5 in the complete
boundary magnitude NAB18. Both $F_{11},F_{22}$ are strictly positive,
since NCE9 has already proved $F_{13},F_{23}\ne0$.

Put $f=(F_{13},F_{23})^{\mathsf T}$, $E=F_{33}$ and retain the exact
diagonal $D_f=\operatorname{diag}(F_{13},F_{23})$. The literal Schur
complement of the third column and its response-coordinate form are
\[
 \begin{split}
 S_F&=\operatorname{diag}(F_{11},F_{22})-ff^*/E,\\
 T_F&=D_f^{-1}S_FD_f^{-*}
  =\begin{pmatrix}
   F_{11}/|F_{13}|^2-1/E&-1/E\\
   -1/E&F_{22}/|F_{23}|^2-1/E
   \end{pmatrix}\succeq0.
 \end{split}                                                     \tag{NCE19}
\]
Indeed $S_F$ is the Gram of the two columns after subtracting their
exact orthogonal components along $v$, proving positivity even if
those two residual columns are dependent. The displayed congruence
proves every entry, including its negative real off-diagonal value.

Define the two scalar observation amplifications without changing any
source or vector by
\[
 \alpha=\frac{E F_{11}}{|F_{13}|^2},\qquad
 \beta=\frac{E F_{22}}{|F_{23}|^2}.
 \tag{NCE20}
\]
The squared norm of the orthogonal projection of $v$ onto the two
mutually orthogonal original columns is
$|F_{13}|^2/F_{11}+|F_{23}|^2/F_{22}$. It cannot exceed $E$:
subtracting that projection leaves an orthogonal square. Consequently
\[
 \alpha,\beta\ge1,\qquad
 \frac1\alpha+\frac1\beta\le1,\qquad
 \alpha\beta\ge\alpha+\beta,
 \qquad
 \boxed{\alpha\beta=
       \frac{\kappa_N^2}{|\mathfrak c_{N,\lambda}|^2}.}
 \tag{NCE21}
\]
The final equality follows by multiplying NCE20 and using NCE9 and
NCE18. This evaluates the product of the two original denominator
amplifications in the complete action magnitude and the actual current.
It does not assign either factor from the other or from its rank.


\begin{thebibliography}{99}
\bibitem{prog:arrival02}
Supplied mathematical continuation, 19 September 2026,
\emph{The original native numerator and the complete retained arithmetic class},
incoming transcript \texttt{Pasted text.txt}, attachment
\texttt{776c17bb-b7af-4c3f-9f05-77299b992322}.
The complete source bytes and precise SHA-256 are retained in
\texttt{SOURCE\_PINS.json}. NCB1--35 verifies its equations (3)--(17)
and supplies the complete minor transitions and action-rank calculation.

\bibitem{prog:DNE}
Split-Zero research programme, \emph{Current kernel and multiplicity proofs},
complete supplied source \texttt{14\_CURRENT\_KERNEL\_AND\_MULTIPLICITY\_PROOFS.tex},
DNE1--15, original 1-based source lines 20772--21116.
Frozen source SHA-256:
\path{8b7b185c31cb7f863dc857bc6978458522002974ddf2a76495fe8de619486cd8}.
The full conductor map retains the original unscaled $E_A$;
the exact companion convention is checked against XIM7 and XIM12
in \texttt{LRC\_TRANSPORT\_COMPLETE.tex}.

\bibitem{prog:FCI}
Split-Zero research programme, 17 September 2026,
\emph{The original comparison cone, its native spectrum, and its arithmetic map},
complete source \texttt{FCI\_SOURCE.tex}, FCI.1--39.
The present morphism uses FCI.33--39 and its complete primary unit;
the original Taylor and theta-primitive provider is FDC.13--15.
No literal identity of the relative conductor denominator with theta
boundaries is asserted by that source or this continuation.

\bibitem{prog:arithmetic}
Split-Zero research programme, 2026,
\emph{Signed arithmetic integration},
\texttt{SIGNED\_ARITHMETIC\_INTEGRATED.tex}, original attained
moment quotient and full arithmetic source;
and \emph{The original comparison cone}, FCI.38--39.
NAB3--19 and NCRP8 prove directly every finite minimum and boundary
identity used in the new calculation.

\bibitem{prog:EC}
Supplied continuation, \emph{The retained eigenclass's marked pairing},
complete EC source in \texttt{02\_Untitled 1775.md},
original 1-based source lines 9401--10815;
EC37--41 start at lines 10731, 10750, 10762, 10775 and 10790.
The exact current and response calculation is re-proved in NCE1--21.
The later parity-dependent analytic radius is not newly certified here.

\bibitem{human:boyd-vandenberghe}
Stephen Boyd and Lieven Vandenberghe,
\emph{Convex Optimization}, Cambridge University Press, 2004.
Appendix A.5.5, p.~650, and Appendix C.4.1, p.~673,
for positive Schur complements and block elimination.
\url{https://web.stanford.edu/~boyd/cvxbook/bv_cvxbook.pdf}.
The required complex Hermitian square identities are proved in the text.
The programme's pinned source records direct inspection on
17 September 2026; a web fetch on 19 September returned an error,
so no new page inspection is claimed for that later fetch.

\bibitem{human:dlmf18}
T.~H. Koornwinder, R. Wong, R. Koekoek, R.~F. Swarttouw and
W.~P. Reinhardt, ``Orthogonal Polynomials,'' Chapter 18 of the
\emph{NIST Digital Library of Mathematical Functions}.
General recurrence: \url{https://dlmf.nist.gov/18.2}, section (iv).
Meixner--Pollaczek weight and norm:
\url{https://dlmf.nist.gov/18.19}, equations (18.19.8)--(18.19.9).
Recurrence: \url{https://dlmf.nist.gov/18.22.E8}.
These sections were consulted on 19 September 2026.
For NCRP25 the exact dictionary is
$p_{n,s}(y)=n!P_n^{(s/4)}(y/2;\pi/2)$ with complete mass
$(2\pi)^{s/2}$; no arithmetic measure is replaced by that classical one.

\bibitem{human:globalization2026}
\emph{A Note on the Globalization Semiring G(Z)}, Zenodo,
12 March 2026. Deposited creators: ChatGPT, 5.4 and Gemini, 2.5 DeepThink.
Originating programme directed by KokunoYumeto with AI assistance,
as expressly identified by the programme author.
\url{https://doi.org/10.5281/zenodo.18976982}.
This attribution does not claim global novelty.

\bibitem{human:splitzero2026}
The Clankers (deposited creator), \emph{A Secondary Note on Split-Zero
Globalization: Semimodule Classification, Multiplicative Reconstruction,
and All-Orders Zeta Deformation}, Zenodo, record publication date
20 March 2026. Originating programme directed by KokunoYumeto
with AI assistance. \url{https://doi.org/10.5281/zenodo.19120905}.
The deposited PDF prints 29 April 2026; record and printed dates
remain distinct. No global priority claim is made.
\end{thebibliography}

\end{document}
